\section{模型评价与改进}

\subsection{模型的优点}

\begin{enumerate}
  \item \textbf{几何边界明确，关键结论可复核。}
  定位区域、保证接收域的包含关系与近似误差均有明确几何含义。

  \item \textbf{四个问题之间形成连续的决策链。}
  模型从交会定位、选点扩展到扫描、路径清除和定向源决策，前后直接复用。

  \item \textbf{理论完成条件与模拟结果相互印证。}
  两类布局和清除点列给出条件，问题三、四的50次日志清除比例均为 $100\%$。
\end{enumerate}

\subsection{模型的局限性}

\begin{enumerate}
  \item \textbf{有限离散与替代目标限制了最优性结论。}
  有限场景、Fisher 指标和有限状态只能支持当前候选集内的比较。

  \item \textbf{搜索、定位与路径规划仍采用分阶段优化。}
  补充定位、清除邻域和访问顺序未统一优化，不能证明任务总时间最短。

  \item \textbf{极端几何和长尾分支的保证具有条件。}
  近共线容差、半径超过 $80\ \mathrm m$ 的截断及有限状态使完成性结论具有条件。
\end{enumerate}

\subsection{模型的改进方向}

\begin{enumerate}
  \item \textbf{采用自适应几何精度。}
  按边界夹角和区域尺度调整外包边数，并稳健处理近共线交点。

  \item \textbf{直接围绕最终评价指标优化选点。}
  在高敏感位置增加场景，直接优化最坏定位半径和行动时间。

  \item \textbf{将定位、清除与路线进行滚动联合规划。}
  将区域收缩、清除邻域和后续访问顺序纳入有限时域模型。
\end{enumerate}

